% ============================================================
% Section 6: Integration with the BX3 Framework
% Recursive Spawning Protocol — Draft
% ============================================================
\section{Integration with the BX3 Framework}
\label{sec:rs-integration}

The Recursive Spawning Protocol is not a module composed alongside the BX3 Framework as a separate component. It is a \emph{structural instantiation} of the BX3 three-layer accountability architecture: every RSP mechanism maps onto exactly one BX3 layer, and each such mapping is \emph{minimal and structurally necessary}. There are no advisory constraints, no optional gates, and no bypass paths for any machine actor under any operational circumstance. The three mappings are: Purpose Layer as the exclusive accountability anchor, Bounds Engine as depth enforcement, and Fact Layer as the forensic ledger. No proper subset of these layers achieves the same guarantees.

% ----------------------------------------------------------
\subsection{Purpose Layer: Accountability Anchor}
\label{sec:rs-purpose-anchor}

The Purpose Layer occupies two exclusive roles in the Recursive Spawning Protocol, and neither can be performed by any other layer.

\medskip
\noindent\textbf{Exclusive origin of spawn authorization.} At root provisioning, the Purpose Layer defines the spawn authorization policy $P_{\mathrm{spawn}}$ and records it in the Fact Layer authorization ledger. $P_{\mathrm{spawn}}$ mandates two conditions for any valid child spawn: (i) $R(n_{i+1}) \subset R(n_i)$ — strict scope refinement — and (ii) operational necessity that cannot be met within the parent's operating scope. No machine actor may issue a spawn warrant in the absence of a matching Purpose Layer authorization record. The Bounds Engine evaluates against $P_{\mathrm{spawn}}$; the Fact Layer enforces it; neither originates it. If the Purpose Layer were not the exclusive origin, a machine actor could define its own spawn authorization policy and begin delegating to descendants with no human oversight — eliminating the traceability of the chain to any human principal.

\medskip
\noindent\textbf{Exclusive terminus of escalation.} When a chain reaches $D_{\mathrm{ESCALATE}}$, the protocol compulsory-escalates to $h(n)$. This terminus is non-collapsing: for any $n_i \in C$, $h(n_i) = h(n_0)$. The human anchor does not migrate down the chain as depth increases. The human issues one of $\{\texttt{ACK},\ \texttt{RESOLVE},\ \texttt{REJECT}\}$; no machine actor may issue any of these. If any machine actor could absorb an exception or issue a terminal directive, the chain would terminate at that actor rather than at a human — defeating the upstream BX3 accountability guarantee. The Purpose Layer's exclusivity here is not a design preference; it is the load-bearing constraint that makes the guarantee structurally operative.

The structural necessity of the Purpose Layer in RSP is therefore established by two exclusive roles that no other layer can occupy: it alone originates spawn authorization policy, and it alone serves as the chain's terminal accountability terminus.

% ----------------------------------------------------------
\subsection{Bounds Engine: Depth Enforcement}
\label{sec:rs-bounds}

The Bounds Engine enforces constrained execution evaluation at every spawn event. It evaluates proposed actions against the operating range defined by the Purpose Layer and enforces the depth constraint that prevents unbounded chain growth. The Bounds Engine cannot exceed its constraints under any operational circumstance.

\medskip
\noindent\textbf{Depth evaluation at spawn.} At each spawn event, the Bounds Engine evaluates current chain depth $|C|$ against $D_{\mathrm{max}}$ (anchored in the Fact Layer authorization ledger). If $|C| \geq D_{\mathrm{max}}$, the spawn is not executed. The Bounds Engine transitions to \texttt{ESCALATING} and initiates the Bailout Protocol. This enforcement is deterministic: no discretion to exceed $D_{\mathrm{max}}$ regardless of urgency or operational exigency. The depth constraint is not a heuristic — it is a hard architectural boundary enforced by the Fact Layer pre-commit hook.

\medskip
\noindent\textbf{Structural necessity.} The structural necessity of the Bounds Engine follows from the drift triad: depth limits alone cannot prevent scope-creep drift (as proven in Section~\ref{sec:rs-drift}). The Bounds Engine enforces the depth constraint; the Fact Layer enforces scope refinement; together they enforce both conjuncts of the drift prevention condition. Neither layer alone is sufficient. A system with only a Bounds Engine (no Fact Layer pre-commit) could allow a machine actor to override depth constraints in exigent circumstances. A system with only a Fact Layer pre-commit hook (no Bounds Engine) would lack an architectural component dedicated to depth evaluation at spawn time — depth constraints would be enforced only implicitly, without a dedicated enforcement point.

The Bounds Engine and the Purpose Layer operate in coordination but are architecturally distinct. The Purpose Layer defines what authorized scope refinement means; the Bounds Engine evaluates whether a specific spawn event violates the depth constraint. Neither can substitute for the other.

% ----------------------------------------------------------
\subsection{Fact Layer: Forensic Ledger}
\label{sec:rs-fact-ledger}

The Fact Layer provides deterministic enforcement and forensic auditability. It maintains the append-only ledger that makes the RSP chain auditable at any future time, and enforces the write protocol that makes RSP tamper-evident.

\medskip
\noindent\textbf{Append-only forensic ledger.} The Fact Layer records every spawn event, every state transition, every Purpose Layer directive, and every unresolved exception. Each entry is timestamped, attributed to the originating node, and hash-chained to its predecessor. Hash-chaining ensures no entry can be inserted, deleted, reordered, or altered after recording: any tampering breaks hash continuity and is detectable by any observer with ledger read access. The ledger is the runtime artifact that distinguishes RSP chains from conventional delegation chains, which exist only as forensic reconstructions attempted after the fact from mutable logs.

\medskip
\noindent\textbf{Immutable parent pointer.} Once $\alpha(n)$ is established at node provisioning, no actor — including the Purpose Layer post-provisioning — may modify it. The Fact Layer write protocol rejects any attempt to overwrite $\alpha(n)$. Chain topology changes occur only through authorized spawn operations recorded in the ledger; no mechanism for retroactive pointer manipulation exists. The parent pointer is not merely access-controlled — it is structurally immutable by protocol definition. This eliminates the re-parenting drift pathway (Failure Mode 3) by making retroactive chain rewriting structurally impossible.

\medskip
\noindent\textbf{Pre-commit hook.} At every spawn event, the Fact Layer pre-commit hook verifies conditions (a) $R(n_{i+1}) \subseteq R(n_i)$ and (b) $\mathrm{depth}(n_{i+1}) \leq D_{\mathrm{max}}$ \emph{before} the spawn is recorded. Rejected spawns are logged; the Bounds Engine may not execute them; \texttt{ESCALATING} is entered if the unmet condition persists. The pre-commit hook operates identically on all machine actors; no privileged actor capable of bypassing it exists. The pre-commit hook makes both the depth bound and the scope refinement \emph{structural invariants} — not policy conventions overridable in exigent circumstances.

The structural necessity of the Fact Layer is established by the immutable-parent-pointer requirement: without the Fact Layer write protocol, immutability is a promise rather than a constraint. A mutable parent pointer enables all four drift failure modes; without the pre-commit hook, scope refinement and depth bounds are policy documentation — enforceable until an actor with sufficient privilege decides they are inconvenient. The Fact Layer is the structural enforcement layer that makes policy effective.

% ----------------------------------------------------------
\subsection{Structural Minimality}
\label{sec:rs-minimality}

The three-layer mapping is minimal-sufficient. Removing any layer causes the corresponding guarantee to fail, and no proper subset of layers achieves the RSP's correctness properties.

\begin{table}[h]
\centering
\begin{tabular}{p{3.0cm} p{5.5cm}}
\toprule
\textbf{Removed layer} & \textbf{Guarantee that fails} \\
\midrule
Purpose Layer        & Exclusive human terminus eliminated; chain can terminate in machine actors; spawn authorization originates in machine actors rather than human intent \\
Bounds Engine       & No architectural component enforces depth constraint at spawn; depth is enforced only implicitly; depth-limit insufficiency theorem applies \\
Fact Layer          & Immutable parent pointer eliminated; retroactive chain manipulation possible; depth bound and scope refinement become overridable conventions \\
\bottomrule
\end{tabular}
\caption{Consequences of layer removal for RSP correctness properties.}
\label{tab:layer-minimality}
\end{table}

No subset of two layers achieves all three correctness properties simultaneously. A system without the Purpose Layer cannot guarantee that spawn authorization traces to human intent. A system without the Bounds Engine cannot guarantee depth-bounded growth. A system without the Fact Layer cannot guarantee chain integrity or forensic auditability. The three-layer architecture is the minimal structure that makes RSP's correctness properties structurally guaranteed rather than conditionally expected.